\section{Introduction}\label{sec:intro}

Quick commerce (q-commerce) promises groceries and household essentials within minutes, fulfilled from small urban warehouses called \emph{dark stores}. India has become its largest laboratory: in the April--June 2026 quarter Blinkit alone reported 2,443 dark stores and 331 million orders \citep{eternal2026q1}, and Swiggy Instamart a further 1,171 stores \citep{swiggy2026q1}. The model rests on a simple geometry: a store that serves only a 1.5--3~km radius can meet a 10-minute promise without riders travelling fast, which is how platform founders defended it \citep{palicha2022,dhindsa2021}.

That defence has been tested. After a nationwide strike by gig workers on 31 December 2025, the Union Labour Ministry intervened and in January 2026 Blinkit, Zepto, Swiggy Instamart and Zomato dropped ``10-minute delivery'' from their branding \citep{outlook2026drop}. In August 2026 the Parliamentary Standing Committee on Commerce recommended guidelines against ``unrealistic speed targets'' set by dispatch algorithms, but did not say what makes a target unrealistic \citep{outlook2026panel}. The policy debate therefore lacks a quantitative answer to a classical OR question: \emph{what does each minute of delivery promise cost, and who pays for it---the platform, through more stores, or the rider, through higher speed?}

This paper answers that question for Bengaluru, India's most contested q-commerce market, and situates it in the OR literature. We pose three research questions:

\begin{description}
  \item[RQ1] How has OR research on on-demand and last-mile delivery evolved, which methods dominate (location, assignment, routing), and where are the gaps for quick commerce and courier welfare?
  \item[RQ2] For Bengaluru, how do the number of dark stores, rider requirements and cost change as the delivery promise relaxes from 10 to 15, 20 and 30 minutes under a safety-capped rider speed?
  \item[RQ3] Which promise level and network design balance coverage, cost and rider safety, and what does this imply for defining a ``realistic'' speed target?
\end{description}

Our contributions are fourfold. First, a bibliometric map of OR research on last-mile and on-demand delivery from 2015 to 2026, which shows a fast-growing field (\BibGrowthPct\% a year) centred on vehicle routing with new delivery modes, in which only \BibQComm{} of \BibIncluded{} articles address quick commerce, none studies India and only \BibWelfare{} consider courier welfare. Second, an integrated and fully reproducible model chain---coverage-constrained facility location, transportation, assignment and PERT---built from public Indian data, together with an audit showing why widely used Kaggle delivery datasets are unsuitable for spatial modelling. Third, a ``price of speed'' curve for Bengaluru. Fourth, evidence for regulators on where a delivery promise stops being achievable through network design and starts depending on rider speed.

Section~\ref{sec:lit} reviews the literature; Section~\ref{sec:method} describes the bibliometric protocol and modelling framework; Section~\ref{sec:data} presents the case and data; Section~\ref{sec:model} formulates the models; Section~\ref{sec:results} reports results; Section~\ref{sec:discussion} discusses implications and limitations; Section~\ref{sec:conclusion} concludes.
